\begin{appendix}
\chapter{Anexos}\label{chap:ann1}

\section{Anexo A: Arquitectura del Sistema HFL}
\label{anx:arquitectura}

La arquitectura desplegada consta de tres capas principales y una extensión FOG opcional:

\begin{itemize}
    \item \textbf{Capa Edge:} 2 $\times$ ESP32-S3 (nodo normal + nodo simulado), ejecutando inferencia TinyML y comunicación MQTT cifrada con ASCON-128.
    \item \textbf{Capa Fog:} Raspberry Pi 4 actuando como broker MQTT, gateway de entrenamiento local y puente HTTP hacia el servidor.
    \item \textbf{Capa Cloud:} 1 $\times$ PC (servidor FastAPI), coordinando la agregación FedAvg y el dashboard analítico.
    \item \textbf{Extensión FOG:} Raspberry Pi líder y peer usando \texttt{fog/weights} y \texttt{fog/global\_model} para agregar entre gateways antes del PC.
\end{itemize}

\section{Anexo B: Archivos del Sistema}
\label{anx:archivos}

\begin{table}[H]
\centering
\caption{Archivos principales del sistema HFL v7.}
\small
\begin{tabular}{l l l}
\toprule
\textbf{Archivo} & \textbf{Dispositivo} & \textbf{Función} \\
\midrule
\texttt{main\_edge\_node\_normal.cpp}    & ESP32-S3 & Nodo de tráfico normal \\
\texttt{main\_edge\_node\_simulated.cpp} & ESP32-S3 & Nodo de tráfico mixto \\
\texttt{ascon128.h}                      & ESP32-S3 & Cifrado ASCON-128 (C++) \\
\texttt{gateway\_hfl.py}                 & Raspberry Pi 4 & Gateway federado \\
\texttt{gateway\_hfl\_fog.py}             & Raspberry Pi 4 & Gateway FOG líder/peer \\
\texttt{ascon128.py}                     & RPi / PC & Cifrado ASCON-128 (Python) \\
\texttt{ascon\_metrics.py}               & RPi / PC & Métricas de cifrado \\
\texttt{plain\_metrics.py}               & RPi / PC & Métricas de baseline sin ASCON \\
\texttt{server\_hfl.py}                  & PC & Servidor central FedAvg \\
\texttt{server\_hfl\_fog.py}              & PC & Servidor para agregación FOG \\
\bottomrule
\end{tabular}
\end{table}

\section{Anexo C: Repositorio del Proyecto}
\label{anx:repositorio}

El código fuente completo, datasets y scripts de entrenamiento están disponibles en el repositorio del proyecto. Las carpetas relevantes para esta revisión son \texttt{hfl\_v7-RN} (MLP/RN con ASCON), \texttt{hfl\_v7-no-ascon-RN} (baseline MLP/RN con JSON plano), \texttt{hfl\_v7-CNN} (modelo CNN-1D con ASCON, estándar y FOG), \texttt{NotebooksAndModels} (entrenamiento offline y exportación de pesos), \texttt{Analisis de Modelos} (post-procesamiento por variante: \texttt{RN}, \texttt{no-ascon}, \texttt{CNN}, \texttt{CNN\_FOG}, \texttt{NIST}) y \texttt{Papers} (literatura base).

\section{Anexo D: Resumen de literatura clave incorporada}
\label{anx:literatura}

La Tabla~\ref{tab:anx_literatura_clave} resume los trabajos más recientes que sustentan las decisiones de diseño y los trabajos futuros del Capítulo~\ref{chap:conc}, distinguiendo lo que aporta cada referencia.

\begin{longtable}{@{}p{0.32\textwidth} p{0.20\textwidth} p{0.42\textwidth}@{}}
\caption{Literatura clave y aporte concreto a la tesis.}
\label{tab:anx_literatura_clave}\\
\toprule
\textbf{Trabajo} & \textbf{Tipo} & \textbf{Aporte a la tesis} \\
\midrule
\endfirsthead
\toprule
\textbf{Trabajo} & \textbf{Tipo} & \textbf{Aporte a la tesis} \\
\midrule
\endhead
Kaur~\textit{et~al.}~2025 \cite{Kaur2025_ASCON_Survey} & Survey ASCON (ACM CSUR) & Justifica elección de ASCON-128 sobre 128a y guía contramedidas SCA. \\
Mirzaali~\textit{et~al.}~2025 \cite{Mirzaali2025_LWC_Permutations} & Análisis crítico LWC & Argumenta no usar permutaciones \emph{ad-hoc}; sugiere ataques replicables. \\
Tran~\textit{et~al.}~2026 \cite{Tran2026_EdgeTrustShard} & FL Bizantino MCU & Encuadra brecha de memoria HFL-MCU y motiva trabajo futuro Bizantino. \\
Ficco~\textit{et~al.}~2024 \cite{Ficco2024_FL_TinyML_Onboard} & FL+TL on-board MCU & Antecedente directo de FL+TinyML on-board en ESP32 y RPi. \\
Dini~\textit{et~al.}~2024 \cite{Dini2024_EmbeddedIIoT_Review} & Survey IIoT/AI embed. & Encuadra \emph{toolchains} TinyML y trade-offs SW/HW. \\
Ni \& Li 2024 \cite{Ni2024_IIoT_Security_Survey} & Survey IIoT seguridad & Sustento del problema (IIoT, ML, retos). \\
Amanullah~\textit{et~al.}~2020 \cite{Amanullah2020_DL_BigData_IoT_Security} & Survey DL+IoT & Antecedente histórico de DL para seguridad IoT. \\
Bagaa~\textit{et~al.}~2020 \cite{Bagaa2020_ML_IoT_Framework} & SDN+ML para IoT & Antecedente de framework ML para seguridad IoT. \\
Hajj~\textit{et~al.}~2023 \cite{Hajj2023_CrossLayerFL_IDS} & FL cross-layer IDS & Comparación directa contra FL para IDS IoT. \\
Chaurasia~\textit{et~al.}~2024 \cite{Chaurasia2024_FL_ResNet_IIoT} & FL+ResNet IIoT & Comparación con cliente FL más expresivo. \\
Roy~\textit{et~al.}~2023 \cite{Roy2023_FL_AdaptedModel} & FL personalizado & Refuerzo del trabajo futuro de personalized FL. \\
Saravana Balaji~\textit{et~al.}~2023 \cite{SaravanaBalaji2023_EdgeIoT_DL} & Edge sanitario+LABE & Antecedente de Edge+cifrado para IoT. \\
Wang~\textit{et~al.}~2019 \cite{Wang2019_DataFusion_CPSS} & Data fusion CPSS & Sustento conceptual del modo FOG. \\
Rani~\textit{et~al.}~2024 \cite{Rani2024_IoT_IDS_Datasets} & Integración datasets IoT-IDS & Sustento de la línea de trabajo futuro sobre datasets reales. \\
\bottomrule
\end{longtable}

\end{appendix}
